En este proyecto se estudia el problema de asignar los viajes de RED a una flota de buses el\'ectricos, determinando simult\'aneamente sus secuencias de operaci\'on y necesidades de recarga, con el objetivo de minimizar los costos operacionales y asegurar que todos los viajes sean cubiertos de manera factible.

Como primera aproximaci\'on, se analizan los datos disponibles de rutas, expediciones, horarios, paraderos, caracter\'isticas de los buses, electroterminales y tarifas el\'ectricas, considerando inicialmente un d\'ia laboral, que concentra 417 rutas y 64.502 expediciones. A partir de la revisi\'on bibliogr\'afica se estudian distintas metodolog\'ias de resoluci\'on, entre ellas Adaptive Large Neighborhood Search, Set Partitioning con Generaci\'on de Columnas y cortes de Benders, y una estrategia de descomposici\'on espaciotemporal. Esta \'ultima se propone como enfoque principal, al dividir el problema primero por electroterminal y posteriormente mediante ventanas horarias superpuestas, reduciendo la escala de los subproblemas a resolver. El presente informe corresponde a una etapa preliminar del proyecto, por lo que se establecen las bases del modelo, los indicadores de desempe\~no y las metodolog\'ias que ser\'an implementadas y evaluadas en las siguientes etapas.

\vfill
\noindent {\bf Palabras Clave}: Electric Vehicle Scheduling Problem, buses el\'ectricos, RED Metropolitana de Movilidad, programaci\'on de recargas, descomposici\'on espaciotemporal, optimizaci\'on de flotas.
